\documentclass[12pt]{article}
\setlength{\oddsidemargin}{27mm}
\setlength{\evensidemargin}{27mm}
\setlength{\hoffset}{-1in}

\setlength{\topmargin}{20mm}
\setlength{\voffset}{-1in}
\setlength{\headheight}{42pt}
\setlength{\headsep}{15pt}


\setlength{\textheight}{215mm}
\setlength{\textwidth}{155mm}


\newcommand\tab[1][1cm]{\hspace*{#1}}
\renewcommand{\thefootnote}{\fnsymbol{footnote}}
\renewcommand{\labelitemi}{$\diamond$}


\usepackage{fancyhdr}
\usepackage[table]{xcolor}
\usepackage{setspace}
\usepackage{graphicx}

\renewcommand{\headrulewidth}{0pt}
\renewcommand{\footrulewidth}{0.5pt}

\definecolor{signaturegray}{gray}{0.98}

\pagestyle{fancy}
\fancyhf{}
\chead{\includegraphics[width=\textwidth]{utadlogo.jpg}}
\rfoot{Page \thepage}
\lfoot{PhD Thesis plan in Electrical and Computer Engineering of UTAD}


\begin{document}

\doublespacing

\baselineskip 12pt

\vspace*{25mm}

\begin{center}

\textbf{\Large PhD Thesis Plan} \\
\vspace{10mm}
\textbf{In Electrical and Computer Engineering}\\
\vspace{20mm}
{\large University of Trás-os-Montes and Alto Douro}\\
\vspace{20mm}

According to Article 14 (g) of the\\
{\scriptsize Regulation of Study Cycles Leading to the UTAD's Doctoral Degree\\
DR, 2.ª série – N.º 133 – Regulamento n.º 656/2016 de 13 de julho de 2016}\\
\end{center}

\vspace{20mm}

{\scriptsize Candidate name:\tab[7cm]Number:}\\
Put your name here \tab[7cm]123456\\
\vspace{3mm}

{\scriptsize Provisional title:}\\
Thesis title here\\
\vspace{3mm}

{\scriptsize Supervisor 1 name:}\\
Supervisor name here\\
\vspace{3mm}

{\scriptsize Supervisor 2 name:}\\
Supervisor 2 name here\\
\vspace{3mm}

{\scriptsize Supervisor 3 name:}\\
Supervisor 3 name here\\
\vspace{3mm}

\pagebreak

\section{Theme/Topic} 
This point should include a description of the current state of knowledge in the proposed domain and properly frame the thesis topic.
Example: One area expected to grow in the scope of Electricity Markets (EM) is Demand Response (DR), as it appears as a very promising opportunity for consumers and brings several advantages for the whole system \cite{Hunt}. This is due to the fact that power systems’ infrastructure is highly capital-intensive and DR is one of the cheaper resources available to prevent investment needs, by peak-shaving and strategic load curtailment in congestion situations \cite{Aalami}. On the other hand, DR programs can provide the system operator with a determined load curtailment capacity which is highly valuable to deal with unexpected changes in both supply and demand levels. References like \cite{Walawalkar} and (sub)sections should be numbered like in this sample (...).

\section{Place of work execution}
The work of this thesis is developed within the scope of INESC Research Group. Alongside these institutions are UTAD - University of Trás-os-Montes and Alto Douro and xxxxxx.

\section{Motivation and contributions}
In this section, the motivation that led to the thesis formulation should be presented, as well as the innovative contributions that this work aims to establish.

\section{Objectives/Goals}
In this section, the main objectives to be achieved with this thesis should be presented. Example:
The main objective of the proposed PhD work is to develop, implement and validate a methodology to support decision making concerning Demand Response (DR) programs and contracts design and use. 
In the path to attain its main objective, the PhD research will also have in mind the following goals:
- Updated state of the art concerning DR
- Design of modified and new DR programs, contracts and business models
- Modeling of DR programs, contracts and business models, according to the state of the art and to the ones proposed in the scope of the PhD work
- Modeling of the players involved in DR, namely consumers, ISOs, VPPs, and CSPs
- Developing and implementing a machine learning (ML) based methodology to support DR programs use. This methodology will be based on clustering and classification techniques, resulting in a rule base concerning DR programs and contracts use. A large set of simulation results will be used as the basis data for the application of the proposed ML method
- Designing and implementing a DR program and contract use evaluation module
- Testing and validating the developed models and the DSS
- Writing at least 6 papers for top level conferences and 3 papers for SCI journals.
(...)


\section{Schedule/Work plan}

\begin{table}
\centering
%\begin{scriptsize}
\caption {Tasks development}
\label{tab:340W}
\begin{tabular}{| c | c | c | c |c | c | c |c | c | c |c | c | c |c | }
\hline 
  month & 0 & 3 & 6 & 9 & 12 & 15 & 18 & 21 & 24 & 27 & 30 & 33 & 36 \\
\hline 
T1  & \cellcolor{gray} & \cellcolor{gray} & \cellcolor{gray} &   &   &   &   &   &   &   &   &   &   \\
\hline
T2  &  &  &  & \cellcolor{gray} & \cellcolor{gray}  & \cellcolor{gray} &   &    &  &  &  &  &  \\
\hline
T3  &  &  &  &  & \cellcolor{gray} & \cellcolor{gray}  & \cellcolor{gray} &  \cellcolor{gray} & \cellcolor{gray}  &  &  &  & \\
\hline
T4  &  &  &  &  &  &  &  &   &   &   &   &   & \\
\hline
T5  &  &  &  &  &  &  &  &   &  &  &  &  & \\
\hline
T6  &  &  &  &  &  &  &  &   &  &  &  &  & \\
\hline
T7  &   &  &  &  &  &  &  &   &  &  &  &  & \\
\hline
\end{tabular}
%\end{scriptsize}
\end{table}


\begin{itemize}
\item \textbf{Task 1:} Example, this task will take place during the first 6 months of the work.  This task will begin with a solid study about the existing DR programs and models. The study of electricity market models is also necessary for better comprehension of the functioning of the integrated DR programs. Several real and relevant electricity markets will be studied, identifying the commercialization and penalization mechanisms used in each one.  The market operation model is of crucial importance for player’s behavior, which should be strategically adapted in function of that model. After this, the problems, disadvantages and weaknesses concerning the existing DR programs will be identified.
\item \textbf{Task 2:};
\item \textbf{Task 3:};
\item \textbf{Task 4:};
\item \textbf{Task 5:};
\item \textbf{Task 6:};
\item \textbf{Task 7:}.
\end{itemize}


\begin{thebibliography}{9}
\bibitem{Hunt}
Sally Hunt, ``Making Competition Work in Electricity,'' John Wiley , pp.23-30, Fev, 2002.

\bibitem{Aalami}
H Aalami, M Moghaddam, G Yousefi, “Demand response modeling considering Interruptible/Curtailable loads and capacity market programs”,  Applied Energy, Vol. 87, no. 1, Elsevier, January 2010.

\bibitem{Walawalkar}
R Walawalkar, S Fernands, N Thakur, K Chevva, “Evolution and current status of demand response (DR) in electricity markets: Insights from PJM and NYISO”, Energy, Vol. 35, no. 4, Elsevier, April 2010.

\end{thebibliography}
\vspace{20mm}
Vila Real, xx of Month of 2021\\[-5mm]


\begin{center}
\textbf{Supervisor 1} \\
\colorbox{signaturegray}{\makebox[0.45\textwidth]{
\raisebox{-35mm}{
\rule{5cm}{0.1pt}}}} \\
Name of Supervisor 1
\vspace{5mm}


\textbf{Supervisor 2} \\
\colorbox{signaturegray}{\makebox[0.45\textwidth]{
\raisebox{-35mm}{
\rule{5cm}{0.1pt}}}} \\
Name of Supervisor 2
\vspace{5mm}

\textbf{Supervisor 3} \\
\colorbox{signaturegray}{\makebox[0.45\textwidth]{
\raisebox{-35mm}{
\rule{5cm}{0.1pt}}}} \\
Name of Supervisor 3
\vspace{5mm}

\textbf{Student} \\
\colorbox{signaturegray}{\makebox[0.45\textwidth]{
\raisebox{-35mm}{
\rule{5cm}{0.1pt}}}} \\
Student Name
\end{center}

\end{document}
